\Problem{Rutas Baratas}
\Statement
Hemos recopilado abundante información sobre las carreteras locales y
alojamientos de una cierta región que queremos visitar. Nuestro plan
es ir de una ciudad $A$ a otra ciudad $B$, gastando la menor cantidad
de dinero posible. Para toda carretera que conecta dos ciudades $u$ y
$v$ sabemos el coste $\omega(u,v)=\omega(v,u)$ de viajar por dicha
carretera (peajes, gasolina, comidas durante el viaje, \ldots).  Cada
vez que viajamos de una ciudad $u$ a una de sus vecinas $v$ debemos
parar en $v$ y hacer noche; sabemos los costes $\omega'(v)$ de
pernoctar para todas las ciudades $v$ (el coste añadido por $A$ y $B$
a nuestra ruta es 0, ya que son los puntos de origen y de destino).
Todos los costes, de vértices y de aristas, son positivos. Por lo
tanto el coste de la ruta
\[
P = [A, v_1, \ldots, v_n, B]
\]
es
\[
  \text{coste}(P) = \omega(A,v_1) + \omega(v_1,v_2) + \ldots + \omega(v_n,B) 
  + \omega'(v_1)+\ldots+\omega'(v_n).
\]

Escribe un programa en C++ que, dados un garfo no dirigido con pesos
positivos en vértices
y en aristas, y dos vértices $A$ y $B$, devuelve el coste de la ruta más barata para ir de
$A$ a $B$, o una indicación de que no existe tal ruta.

\Input
Todos los datos de entrada son enteros positivos.
La entrada comienza con dos enteros
$2{\le}n{\le}10000$ y $m$, $0{\le}m\le{20n}$.
A continuación, viene una secuencia de $n$ enteros positivos
$\omega'(0), \ldots, \omega'(n-1)$, los pesos $\omega'(u)$ de los $n$ vértices del grafo.
Luego viene una secuencia con las $m$ aristas del grafo en forma de tripletas
${\langle}u,v,\omega(u,v)\rangle$. Los vértices $u$ y $v$ son enteros en el rango
$\{0,\ldots,n-1\}$ y los pesos $\omega(u,v)$ son enteros positivos.
Ningún peso, ni de aristas ni de vértices, es mayor que 100000.
Puede asumirse que no hay aristas paralelas diferentes uniendo un mismo par de vértices
y que no hay ninguna arista que une a un vértice consigo mismo.
Finalmente, la entrada contiene una secuencia de pares
${\langle}A_i, B_i\rangle$, donde los $A_i$'s y los $B_i$'s
denotan vértices del grafo ($0{\le}A_i,B_i<n$).

\Output
Para cada par ${\langle}A_i, B_i\rangle$ de la entrada, el programa escribe el coste $\delta$
de la ruta más barata entre $A_i$ y $B_i$ con el formato \texttt{c($A_i$,$B_i$) = $\delta$}.
% and the cheapest path itself
% (a sequence of vertices starting
% with $A_i$ and ending with $B_i$)
Si no hay rutas entre $A_i$ y $B_i$ el programa escribe
\texttt{c($A_i$,$B_i$) = +oo}. Cada línea de la salida termina
con un salto de línea (\texttt{endl}).

\SampleTwoCol
